\section{Setup, compute and reproduction}\label{app:setup}
\paragraph{Models and training.} U-Net from \texttt{torchcfm} (\texttt{num\_channels=32}, \texttt{num\_res\_blocks=1}, 1,075,361 parameters); Adam ($\mathrm{lr}=10^{-3}$, default betas), batch 128, 50 epochs (468 steps each), no EMA, no learning-rate schedule. Pixels scaled to $[-1,1]$. OT-CFM uses exact OT (POT) on the squared Euclidean cost within each batch. DDPM: linear $\beta$ from $10^{-4}$ to $0.02$, $T=1000$, $\epsilon$-prediction, sampled with deterministic DDIM. Variants change exactly one setting: \texttt{--ot-batch} (OT plan on 512 or 2,048 samples, then split into batches of 128), \texttt{--t-dist} (logit-normal or $\kappa$-mixture), \texttt{--reflow} (60,000 dopri5 pairs from the seed-0 model, fine-tuned from its weights with a new Adam optimiser).
\paragraph{Sampling.} Euler (one network call per step), Heun (two) and classical RK4 (four) on the given time grid; dopri5 from \texttt{torchdiffeq} with $\mathrm{rtol}=\mathrm{atol}=10^{-5}$. Grids: uniform; the shift family $t'=st/(1+(s-1)t)$ with $s\in\{0.25,0.5,2,4\}$; and the $\sqrt\kappa$ grid $t_i=F^{-1}(i/N)$ with $F(t)=\int_0^t\sqrt\kappa/\!\int_0^1\sqrt\kappa$, where $\kappa$ is estimated by finite differences of $v$ on 101 times along 256 dopri5 trajectories (Appendix~\ref{app:d5}). Heun* replaces Heun's last step by an Euler step, RK4* replaces RK4's last step by a Heun step whose end evaluation is at $1-h/2$, so that neither calls the network at $t=1$. Minibatch OT savings were measured for OT batches from 16 to 4,096.
\paragraph{Evaluation.} FID with clean-fid (\texttt{mode=clean}) on 10,000 samples against the 60,000 training images, noise seed 100 (101--105 for the noise-error study); real-versus-real FID is 0.54. % results/e4/mnist_eval.json
Class coverage: $k$-means (scikit-learn, \citealp{pedregosa2011sklearn}) with $k=10$ and $k=20$ as a check, effect size $w=\sqrt{\chi^2/n}$ \citep{cohen1988power}, bootstrap intervals \citep{efron1979bootstrap} from 1,000 resamples (Appendix~\ref{app:course}). Straightness $R$ on 256 paths with 100 Euler steps, noise seed 201. NFE@20 is interpolated linearly in $\log\nfe$. With three seeds the smallest one-sided sign-test $p$-value is $0.5^3=0.125$, which is why we report agreement of signs rather than significance.
\paragraph{One-step samples.} One-step I-CFM samples are 0.045 (RMS per pixel) from the MNIST mean image and 0.39 from the nearest training image, as \eqref{eq:postmean} predicts (Appendix~\ref{app:d3}). % results/e4/mnist_eval.json one_step
\paragraph{Diffusion as CFM.} The score, noise prediction and velocity of a Gaussian path with independent pairs are linear functions of each other via Tweedie's formula \citep{efron2011tweedie}; the identity between one DDIM step and one Euler step on the $\tau$ grid (Proposition~\ref{prop:d9}) was checked numerically to $10^{-14}$.
\paragraph{Compute.} One consumer GPU. All GPU jobs ran one at a time and throttled to about half duty cycle (\texttt{code/common/throttle.py}) so that the machine stayed usable; wall times below are throttled. A base model takes 33--35 minutes unthrottled (seed 0) and 70--72 minutes throttled (seed 2); an FID evaluation over seven Euler NFEs about 20 minutes; dopri5 FID about 17 minutes; the core solver-by-grid study on one model about 1.5 hours. Training wall times: base CFM models 5.3\,h in total (six runs) and DDPM 1.5\,h (two runs); OT batch 512 1.56\,h and 2,048 2.37\,h (the latter resumed after a power loss at epoch 21 from the epoch-20 checkpoint; the time adds both parts); logit-normal and $\kappa$-weighted $t$ 1.42 and 1.29\,h; reflow pair generation 2.05\,h (I-CFM) and 1.84\,h (OT-CFM) plus retraining 1.08 and 1.56\,h. In total, training and pair generation took about 20\,GPU-hours, not counting evaluation. Throttled wall time is a noisy cost measure: the two reflow retrainings did identical work but took 1.08 and 1.56\,h, so the time ratios in Table~\ref{tab:training} are accurate only to tens of percent. % results/e2/*_train.json seconds; results/c5_reflow/pairs_*_s0.json gpu_seconds; results/c3_otbatch/logs/train_otcfm_s0_otb2048.log
\paragraph{Reproduction.}\sloppy\ \texttt{code/README.md} lists the command for every experiment; \texttt{code/run\_queue.sh} runs the final-report experiments in order and resumes after interruption. \texttt{python report/make\_figs.py} and \texttt{python report/make\_tables.py} regenerate every figure and table of this report from \texttt{results/}, and \texttt{report/materials/numbers.json} lists every value they use.
